\section{Определение монотонности функции. Определение локальных минимумов, максимумов, экстремумов. Необходимое условие локального экстремума дифференцируемой функции. Алгоритм нахождение наибольшего наименьшего значения дифференцируемой на отрезке функции.}

\subsection{Определение монотонности функции.}

\begin{definition}[Определение неубывающей функции.]
\[
\forall x_1,x_2\in I:\ (x_1<x_2)\ \Rightarrow\ f(x_1)\le f(x_2).
\]
\end{definition}

\begin{definition}[Определение возрастающей функции.]
\[
\forall x_1,x_2\in I:\ (x_1<x_2)\ \Rightarrow\ f(x_1)< f(x_2).
\]
\end{definition}

\begin{definition}[Определение невозрастающей функции.]
\[
\forall x_1,x_2\in I:\ (x_1<x_2)\ \Rightarrow\ f(x_1)\ge f(x_2).
\]
\end{definition}

\begin{definition}[Определение убывающей функции.]
\[
\forall x_1,x_2\in I:\ (x_1<x_2)\ \Rightarrow\ f(x_1)> f(x_2).
\]
\end{definition}




\subsection{Определение локального минимума.}
\begin{definition}
$x_0$ — точка локального минимума функции $f(x)$, если
\[
\exists\,\delta_0>0\ \forall x\in U_{\delta_0}(x_0):\quad f(x_0)\le f(x).
\]
$f(x_0)$ — локальный min.
\[
\]
\\
$x_0$ — точка строгого локального минимума функции $f(x)$, если
\[
\exists\,\delta_0>0\ \forall x\in \dot U_{\delta_0}(x_0):\quad f(x_0)< f(x).
\]
$f(x_0)$ — строгий локальный min.
\end{definition}

\subsection{Определение локального максимум.}
\begin{definition}
$x_0$ — точка локального максимума функции $f(x)$, если
\[
\exists\,\delta_0>0\ \forall x\in U_{\delta_0}(x_0):\quad f(x_0)\ge f(x).
\]
$f(x_0)$ — локальный max.
\[
\]
\\
$x_0$ — точка строгого локального максимума функции $f(x)$, если
\[
\exists\,\delta_0>0\ \forall x\in \dot U_{\delta_0}(x_0):\quad f(x_0)> f(x).
\]

$f(x_0)$ - строгий локальный max.
\end{definition}

\subsection{Определение локального экстремума.}
\begin{definition}
Точками локального экстремума называются точки min и точки max.
\\
Экстремумы - min и max.
\end{definition}

\subsection{Необходимое условие экстремума дифференцируемой функции}
\begin{theorem}[Теорема Ферма. Необходимое условие экстремума дифференцируемой функции.]
Если $x_0$ — т.\ экстремума функции $f$, то\\
либо $\nexists\, f'(x_0)$,\\ либо $f'(x_0)=0$.
\end{theorem}

\begin{Proof}
Пусть $x_0$ — т. минимума и $\exists\, f'(x_0)$.
Тогда при $x>x_0$ имеем $f(x)-f(x_0)\ge 0$ и $x-x_0>0$, значит
\[
\lim_{x\to x_0+}\frac{f(x)-f(x_0)}{x-x_0}\ge 0.
\]
А при $x<x_0$ имеем $f(x)-f(x_0)\ge 0$ и $x-x_0<0$, значит
\[
\lim_{x\to x_0-}\frac{f(x)-f(x_0)}{x-x_0}\le 0.
\]
Так как $\exists\, f'(x_0)$, эти односторонние пределы равны между собой, поэтому
\[
f'(x_0)=0.
\]

\end{Proof}

\subsection{Алгоритм нахождения наибольшего/наименьшего значения дифференцируемой функции на отрезке}
\begin{lemma}[Алгоритм нахождения наибольшего/наименьшего значения дифференцируемой функции на отрезке]
\[
y=f(x)\ \ \text{непрерывна на }[a,b].
\]

\[
M=\sup_{[a,b]} f(x), \qquad m=\inf_{[a,b]} f(x).
\]

\[
\exists\, x_1\in[a,b]:\ f(x_1)=M,\qquad \exists\, x_2\in[a,b]:\ f(x_2)=m.
\]

\[
x_1,\ x_2 \ \text{среди } a,\ b,\ z_1,\dots,z_n,
\]
где $z_i$ — \textit{критические точки}:
\[
f'(z_i)=0 \ \ \text{или}\ \ \nexists\, f'(z_i).
\]

\textbf{Вывод:} перебираем значения
\[
f(a),\ f(b),\ f(z_1),\dots,f(z_n),
\]
выбираем среди них $\min$ и $\max$.

Этот способ будет единственным возможным, как только мы повысим размерность.
\end{lemma}